\section{Дахин туршилт}
\subsection{Орчин ба автомат тестийн хамрах хүрээ}
БД1-ийн орчин, хугацаа нь II-4-10 багийн шүүмжийн бүртгэлээс, БД2-ын дүн нь IV-2-13 багийн терминалын ажиллуулалтаас авсан. Хугацааг тухайн удаагийн тестийн ажиллуулалтыг танихын тулд үзүүлэв. Хоёр хэмжилтийг өөр орчинд, өөр тооны тестээр хийсэн.
\tbl{Туршилтын хоёр орчин ба бүртгэсэн дүн}{tab:env}{
\begin{tabularx}{\linewidth}{@{}L{22mm}YL{36mm}L{25mm}@{}}\toprule
\hd Хувилбар & \textbf{Орчин} & \textbf{Шалгасан тал} & \textbf{Тестийн дүн}\\\midrule
БД1 & Windows 11, Python 3.14.7, Chrome & II-4-10, шүүмж & 64 / 64, 9.05 с\\
БД2 & Debian 13, Python 3.13.5, pytest 9.1.1, Firefox & IV-2-13, \code{6d9c4a3} & 75 / 75, 5.29 с\\\bottomrule
\end{tabularx}}

Регрессийн суурь 64 шалгалт нь цөм (15), CLI (24), CLI-ийн төлөв хадгалалт (7), сервер (13), Banker-гүй горим (5) гэсэн бүлгүүдтэй. Энэ суурийг хадгалж, хүснэгт~\ref{tab:newtests}-д жагсаасан 11 тохиолдлыг нэмсэн. Хөтчөөр хийсэн товчны болон дүрслэлийн гар шалгалтуудыг 75-д оруулаагүй.
\tbl{Шүүмжийн мөрөөр нэмсэн 11 автомат тест}{tab:newtests}{
\begin{tabularx}{\linewidth}{@{}L{48mm}cY@{}}\toprule
\hd Тестийн файл & \textbf{Тоо} & \textbf{Амжилттай гэж үзэх нөхцөл}\\\midrule
\code{test_banker_off.py} & 1 & Номын горим сэргэнэ. P1-д олгож, P0-д татгалзана. Rollback-ийн дараах төлөв өмнөхтэй ижил.\\\midrule
\code{test_ui_scenarios.py} & 2 & 3×2 жишээнд Release-ийн дараа SAFE. 10×5 жишээнд аюулгүй дараалал бүх процессыг багтаана.\\\midrule
\code{test_cli_commands_file.py} & 3 & Файл ба stdin-ийн гаралт адил. Байхгүй файлд код 2. cp1251 тохиргоонд дамжуулсан гаралтыг UTF-8-аар зөв уншина.\\\midrule
\code{test_web_limits.py} & 3 & 20×8 хүснэгт зөвшөөрөгдөнө. 21×8 ба 20×9 хүснэгтийг татгалзахад хуучин төлөв хадгалагдана.\\\midrule
\code{test_post_headers.py} & 2 & Тохироогүй Content-Type ба Origin-той хүсэлтэд тус бүр 415, 403 ирнэ. Төлөв өөрчлөх функц дуудагдахгүй.\\\bottomrule
\end{tabularx}}

БД1-ийн эцсийн \code{f692fdc} хувилбарыг ажлын салбарт нэгтгэсэн \code{92bc3eb} commit-ийн үед 69 тест тэнцсэн. Дараагийн зургаан шалгалт нь CLI гаралтын кодчилол, вэб хүснэгтийн гурван зааг, HTTP толгойн хоёр тохиолдлыг хамарсан.
\newpage
\subsection{Терминал ба хөтчөөр давтан шалгах заавар}
\lead{1. Тестийн орчин.} Төслийн үндсэн хавтсанд, виртуал орчин идэвхтэй үед дараах командуудыг ажиллуулна. Төслийг анх задалсан бол хоёр дахь командаас өмнө тестийн хамаарлыг эхний командаар суулгана.
\begin{Verbatim}[fontsize=\small,xleftmargin=5mm]
python -m pip install -r requirements.txt
python -m pytest -q
\end{Verbatim}
Шалгасан \code{6d9c4a3} хувилбарын сүүлийн мөр \code{75 passed in 5.29s} байсан. Дахин ажиллуулахад хугацаа өөрчлөгдөж болно.

\lead{2. Номын командын файл.} Дараах командыг терминалд \textbf{нэг мөрөөр} оруулна. Сонголтын ард команд бичсэн файлаа өгснөөр PowerShell, Bash хоёрт ижил аргумент хэрэглэнэ.
\begin{flushleft}\small\ttfamily
python src/cli.py 10 5 7 --max examples/book\_max.txt --allocation examples/book\_allocation.txt --commands examples/book\_commands.txt
\end{flushleft}
Гаралтад Max, Allocation, Need хүснэгтүүдийн хамт дараах шийдвэрүүд ижил дарааллаар гарна. Доор багийн терминалын гаралтаас тухайн мөрүүдийг авлаа.
\begin{Verbatim}[fontsize=\small,xleftmargin=3mm]
Төлөв: SAFE  <P1, P3, P0, P2, P4>
granted: P1-д олголоо; төлөв safe хэвээр
wait: P4 хүлээнэ: R0 хүсэлт=3, Available=2
unsafe: P0-ийн хүсэлтийг олговол unsafe болно; татгалзлаа
Төлөв: SAFE  <P1, P3, P0, P2, P4>
\end{Verbatim}
Гараар оруулах бол \code{--commands} болон түүний файлын замыг орхино. \code{RQ 1 1 0 2} нь P1-ийн хүсэлт, \code{RL 1 1 0 0} нь P1-ээс R0-ийн нэг нэгжийг буцаах команд. Одоогийн төлөвийг \code{*}, гарахдаа \code{exit} гэж бичнэ.

\lead{3. Хөтчийн богино шалгалт.} \code{python src/server.py} гэж ажиллуулаад \url{http://127.0.0.1:5000} хаягийг нээнэ.
\begin{enumerate}
\item «Banker-гүй» жишээг сонгоод «Номын жишээг оруулах» товч дарна. Banker-ийн чагт асаж, Available $(3,3,2)$ болно. Хүснэгт~\ref{tab:req}-ийн гурван хүсэлтийг дараалуулж үр дүнг харьцуулна.
\item Номын жишээг дахин ачаалаад зүүн талын Available-ийн эхний нүдийг 3-аас 4 болгоно. Request, Release идэвхгүй байгааг үзээд «Шалгах» дарна. Батлагдсан Available $(4,3,2)$ болоход товчнууд нээгдэнэ.
\item «Том туршилт 10×5» дээр «+ Нөөц»-ийг гурван удаа дарна. Нөөцийн тоо 8 болоход нэмэх товч идэвхгүй, нэгийг хасахад дахин идэвхтэй болно. Дуусаад серверийг терминалд Ctrl+C дарж зогсооно.
\end{enumerate}
\newpage
\section{Эцсийн үр дүн}
Хувилбаруудыг хэрэглэгчийн ижил үйлдлээр харьцуулахын тулд хүснэгтийн мөр бүрд нэг туршилт сонгов. БД1-ийн баганад шүүмжлэгчийн ажиглалт, БД2-ын баганад засварын дараах ажиглалт эсвэл тестийн үр дүнг оруулав.

\tbl{Ижил үйлдлийг засварын өмнө ба дараа харьцуулсан нь}{tab:change}{
\begin{tabularx}{\linewidth}{@{}L{35mm}YY@{}}\toprule
\hd Туршилт & \textbf{БД1-ийн ажиглалт} & \textbf{БД2-ын үр дүн}\\\midrule
Номын жишээ рүү товчоор буцах & Banker унтраалттай үлдэж, дараалсан хоёр дахь хүсэлт P0-д олгогдов. & Banker ассан. P1-д олгож, P0-д татгалзав. Available $(2,3,0)$ хэвээр.\\\midrule
Батлаагүй засвар & Хүснэгт шинээр харагдсан ч хуучин төлвөөр Request явуулж болдог. & «Шалгах» хүртэл Request, Release түгжигдэж, батлагдсаны дараа шинэ тоогоор ажиллав.\\\midrule
Номын CLI жишээ & Windows дээр зааврын командаас ParserError гарсан. Шүүмжлэгч файлыг дамжуулж байж ажиллуулав. & README-д \code{--commands} ашиглав. Debian дээр файлын болон stdin-ийн гаралт байт бүрээрээ таарав.\\\midrule
Тест ба бэлэн жишээ & 64 тест. 3×2 жишээг шүүмжлэгч гараар тулгасан. & 75 тест. Хоёр бэлэн жишээний төлвийг UI-д өгдөг матрицаас автоматаар шалгана.\\\bottomrule
\end{tabularx}}

Харьцуулалтаас харахад номын жишээг хэрхэн ачаалснаас хамаарч өөр хариу гардаг байсан зөрүү арилжээ. Файлаар команд унших шинэ замын хувьд хуучин stdin горимтой ижил гаралт өгдгийг шалгаснаар өмнөх хэрэглээг мөн хадгалсан.

\section{Хязгаарлалт}
\lead{Загварын нөхцөл.} Процесс бүрийн Max урьдчилан мэдэгдсэн, процесс дуусахад эзэмшиж буй нөөцөө чөлөөлнө гэж үзнэ. SAFE гэдэг нь эдгээр нөхцөлд бүх процессыг дуусгах боломжтой дараалал олдсоныг илэрхийлнэ. Симулятор үйлдлийн системийн бодит процессуудыг ажиллуулах, бодитоор хүлээж буй хүсэлтийн дарааллаас deadlock илрүүлэх үүрэггүй.

\lead{Төлөв хадгалалт.} Серверийн нэг төлөв санах ойд байна. Ижил серверт холбогдсон табууд тэр төлвийг хуваалцдаг ч таб бүр өөр түүхтэй. Хөтчийг шинэчлэхэд серверийн төлөв үлддэг, серверийг дахин асаахад номын эхний жишээ ачаалагдана. 20×8 нь сургалтын вэб хүснэгтийн ашиглалтын хязгаар юм. Гүйцэтгэлийн дээд чадлыг тогтоох тусгай хэмжилт хийгээгүй.

\lead{Давтан туршсан орчин.} Эцсийн 75 тестийг Debian дээр ажиллуулсан. cp1251 тохиргооны шалгалтыг мөн Debian-д дэд процессоор хийсэн. БД2-ын энэ хувилбарыг Windows PowerShell, macOS дээр тусад нь ажиллуулсан дүн бүртгэгдээгүй. Гарын авлагын 2.2-т туршилтын орчныг хувилбар тус бүрээр бүртгэсэн. Эцсийн release-ийн шошгыг нийтлэх үед туршсан хувилбартай тулгана.
\newpage
\section{Дүгнэлт}
Номын нэг жишээн дээр гурван өөр шийдвэрийн учрыг тоогоор харуулж чадсан. P1 нөөц авахад аюулгүй дараалал үлдэнэ. P4-ийн хувьд яг одоогийн сул R0 хүрэхгүй. P0-ийн хувьд хүссэн нэгж хүрсэн ч түр олгосны дараа эхэлж дуусгах процесс олдохгүй. Сүүлийн хоёр шийдвэрийн шалтгаан өөр боловч аль алинд нь батлагдсан хуваарилалт хадгалагдана.

Шүүмжээр илэрсэн гол сургамж нь цөмийн зөв хариуг хэрэглэгч зөв горим, зөв өгөгдлөөр авч чадаж байгаа эсэхийг тусад нь шалгах шаардлага байв. Жишээ солих хоёр замыг нэгтгэснээр товчоор ачаалсан номын жишээ ч CLI-тэй ижил шийдвэр гаргадаг болсон. Бэлэн жишээний өгөгдлийг тесттэй хуваалцсанаар цаашид хүснэгт өөрчлөгдөхөд аюулгүй дарааллын талаарх дүгнэлтийг дахин шалгах боломжтой.

\subsection*{Багийн оролцоо}
Гишүүдийн ажлыг эх кодын түүх болон БД1-ийн тайлантай тулгав. Сүүлийн баганаас тухайн ажлыг commit-оор эсвэл багийн бичгээр тэмдэглэсэн үүргээр баталгаажуулсныг харж болно.

\tbl{Багийн гишүүдийн ажил ба эх баримт}{tab:team}{
\begin{tabularx}{\linewidth}{@{}L{29mm}YL{42mm}@{}}\toprule
\hd Гишүүн & \textbf{Бүртгэгдсэн ажил} & \textbf{Нотолгоо}\\\midrule
Г.Лувсанжал & Банкны алгоритмын цөм ба тест. & Git \gitlink{120eacd}\\
У.Анужин & Нөөцийн граф ба UI дүрслэл. & Git \gitlink{2bf9eac}\\
Г.Ганжаргал & CLI төлөвийн тест. & Git \gitlink{f4f7950}\\
Б.Лхагвадорж & Вэб хүснэгтийн анхны загварт оролцсон. & БД1 тайлан, хавсралт А\\
Я.Төгөлдөр & Дараалал ба төлөвийн логик шалгасан. & БД1 тайлан, хавсралт А\\
Г.Ганхүү & Баримт бичиг, README боловсруулахад оролцсон. & БД1 тайлан, хавсралт А\\\bottomrule
\end{tabularx}}

Эхний гурван мөрийн ажлыг Git түүхээс шууд шалгаж болно. Үлдсэн гурван үүргийг БД1-ийн хавсралтад бүртгэсэн тул багийн гишүүний нэртэй холбогдох тусгай commit дугаар заагаагүй.
\newpage
\section*{Ашигласан материал}
\addcontentsline{toc}{section}{Ашигласан материал}
Нөөц хуваарилалт, аюулгүй дарааллын тооцоонд сурах бичгийн 8-р бүлгийг, түгжрэлийн ойлголтод Coffman болон хамтран зохиогчдын тойм өгүүллийг ашиглав.
\begingroup
\renewcommand{\refname}{}
\begin{thebibliography}{9}
\bibitem{osc} A. Silberschatz, P. B. Galvin, G. Gagne. \textit{Operating System Concepts}. 10th ed. Wiley, 2018. Chapter 8, “Deadlocks”, §8.6.3. \href{https://www.wiley.com/en-us/operating-system-concepts-10th-edition-p-9781119320913}{Хэвлэлийн газрын бүртгэл}.
\bibitem{coffman} E. G. Coffman, Jr., M. J. Elphick, A. Shoshani. “System Deadlocks”. \textit{ACM Computing Surveys}, 3(2), 1971, pp. 67–78. \href{https://doi.org/10.1145/356586.356588}{DOI 10.1145/356586.356588}.
\end{thebibliography}
\endgroup
